\chapter{Dizajn sistema}

\section{UML dijagrami}

Dizajn sistema je modeliran korišćenjem UML dijagrama koji prikazuju različite aspekte arhitekture.

\subsection{Dijagram klasa}

Dijagram klasa prikazuje strukturnu dekompoziciju sistema i odnose između entiteta.

\begin{figure}[H]
    \centering
    \includegraphics[width=\textwidth]{02-class-diagram.png}
    \caption{Dijagram klasa sistema za udomljavanje životinja}
    \label{fig:class-diagram}
\end{figure}

\subsubsection{Hijerarhija korisnika}

Sistem definiše apstraktnu klasu \texttt{Korisnik} iz koje nasleđuju specifični tipovi korisnika:

\begin{itemize}
    \item \textbf{SistemskiAdmin}: Ima pristup globalnim podacima, može registrovati nova udruženja
    \item \textbf{AdminUdruženja}: Upravlja specifičnim udruženjem, dozvolama i izveštajima
    \item \textbf{Volonter}: Dodaje životinje, upravlja zahtevima, povezuje udomitelje
    \item \textbf{Udomitelj}: Može biti registrovan ili neregistrovan, unosi podatke pri zahtevu
    \item \textbf{NeregistrovaniKorisnik}: Može pregledati životinje i poslati zahtev
\end{itemize}

\subsubsection{Entiteti}

\paragraph{Udruženje}
Reprezentuje organizaciju koja upravlja udomljavanjem. Sadrži:


\paragraph{Životinja}
Reprezentuje životnju dostupnu za udomljavanje. Ima tri stanja:
\begin{itemize}
    \item DOSTUPNA - životnja nije u procesu udomljavanje
    \item ZAUZETA - postoji aktivna prijava za udomljavanje
    \item UDOMLJENA - životnja je udomljena
\end{itemize}

\paragraph{Udomljavanje}
Predstavlja relaciju između životnje i udomitelja sa sledećim atributima:
\begin{itemize}
    \item Napomena za udomitelja
    \item Datum udomljavanje
    \item Status komunikacije (SMS, email, ručno)
\end{itemize}

\paragraph{Donacija}
Donacija može biti:
\begin{itemize}
    \item Opšta - donacija za udruženje (životinja = NULL)
    \item Namenjena - donacija za specifičnu životnju
\end{itemize}

Vrste donacija: hrana, voda, lekovi, oprema, novac.

\paragraph{Izveštaj}
Sistem generiše izveštaje na dnevnom, nedelnom i mesečnom nivou sa sledećom statistikom:
\begin{itemize}
    \item Broj udomljenih životinja
    \item Broj vraćenih životinja
    \item Ukupna donacija po vrsti
\end{itemize}

\section{Principi dizajna}

\subsection{Separacija gladanja (Separation of Concerns)}

Arhitektura je organizovana u tri sloja:

\begin{itemize}
    \item \textbf{Presentation Layer} - Avalonia UI sa MVVM arhitekturom
    \item \textbf{Business Logic Layer} - Poslovne pravile i operacije
    \item \textbf{Data Access Layer} - Komunikacija sa bazom podataka
\end{itemize}

\subsection{SOLID principi}

\begin{description}
    \item[S - Single Responsibility] Svaka klasa ima jednu odgovornost
    \item[O - Open/Closed] Klase su otvorene za proširenje, zatvorene za izmenu
    \item[L - Liskov Substitution] Izvedene klase mogu zamenjivati bazne klase
    \item[I - Interface Segregation] Interfejsi su specifični i maleni
    \item[D - Dependency Inversion] Zavisnosti od abstraktnih, a ne konkretnih klasa
\end{description}

\section{Baza podataka}

Sistem koristi PostgreSQL za trajno čuvanje podataka. Skema je normalizovana na treći normalni oblik (3NF).

\subsection{Ključne tabele}

\begin{itemize}
    \item \texttt{Korisnici} - sve vrste korisnika
    \item \texttt{Udruženja} - registrovana udruženja
    \item \texttt{Životinje} - dostupne životnje
    \item \texttt{Zahtevi} - zahtevi za udomljavanje
    \item \texttt{Udomljavanje} - istorija udomljavanja
    \item \texttt{Donacije} - evidencija donacija
    \item \texttt{Izveštaji} - generirani izveštaji
\end{itemize}

\section{Tehnološki stek}

\begin{table}[H]
    \centering
    \begin{tabularx}{\textwidth}{|l|l|}
        \hline
        \textbf{Komponenta} & \textbf{Tehnologija} \\
        \hline
        Programski jezik & C\# 12 \\
        \hline
        Framework & .NET 10 \\
        \hline
        UI Framework & Avalonia 11.3 \\
        \hline
        Baza podataka & PostgreSQL \\
        \hline
        Verzijski sistem & Git/GitHub \\
        \hline
    \end{tabularx}
    \caption{Tehnološki stek}
    \label{tab:tech-stack}
\end{table}
